\medskip

On considère la fonction $f$ définie sur $\R$ par:
\[f(x) = \dfrac{1}{1 + \text{e}^{- 3x}}.\]

\begin{center}
\psset{xunit=1.25cm,yunit=2.5cm,arrowsize=2pt 3}
\begin{pspicture*}(-2.5,-0.5)(3,1.5)
\psgrid[subgriddiv=4](-3,-1)(3,2)
\psaxes[linewidth=1.25pt,Dx=10,Dy=10]{->}(0,0)(-2.5,-0.5)(3,1.5)
\psaxes[linewidth=1.25pt,Dx=10,Dy=10](0,0)(1.1,1.1)
\psplot[plotpoints=2000,linewidth=1.25pt,linecolor=blue]{-2.5}{3}{2.71828 x 3 mul exp 2.71828 x 3 mul exp 1 add div}
\psplot[plotpoints=2000,linewidth=1.25pt]{-2.5}{3}{3 x mul 4 div 0.5 add}
\uput[l](0,0.5){A}\uput[dr](1,1.25){B}
\uput[dr](1,0){1}\uput[ul](0,1){1}
\psline[ArrowInside=->,linewidth=1.25pt,linecolor=red](0,0.5)(1,0.5)(1,1.25)
\uput[d](0.5,0.5){1}\uput[r](1,0.875){$\frac34$}
\psdots[dotstyle=+,dotangle=45,dotscale=1.5](0,0.5)(1,1.25)
\uput[d](2.8,0.9){\blue $\mathcal{C}_f$}\uput[r](-1.2,-0.4){$\mathcal{T}$}
\end{pspicture*}
\end{center}

\textbf{Partie A : lectures graphiques}

\medskip

\begin{enumerate}
\item 
Comme indiqué sur la figure pour aller de A en B points de la tangente : + 1 puis $+ \dfrac34$, donc coefficient directeur de $\dfrac34$ et l'ordonnée à l'origine est égale à $\dfrac12$.

$M(x~;~y) \in \mathcal{T} \iff y = \dfrac34x + \dfrac12$.
\item 
La fonction semble être convexe sur $]- \infty~;~0]$ et concave sur $[0~;~+ \infty[$.
\end{enumerate}

\bigskip

\textbf{Partie B : étude de la fonction}

\medskip

\begin{enumerate}
\item 


En posant $u(x) = 1 + \text{e}^{- 3x}$, on a pour tout réel $u'(x) = -3\text{e}^{- 3x}$.

De $f(x) = \dfrac{1}{u(x)}$, on a donc $f'(x) = - \dfrac{u'(x)}{u(x)^2} = - \dfrac{-3\text{e}^{- 3x}}{\left(1 + \text{e}^{- 3x}\right)^2} = \dfrac{3\text{e}^{- 3x}}{\left(1 + \text{e}^{- 3x}\right)^2}$. 
\item 
On a pour tout $x \in \R$, \: $\text{e}^{- 3x} > 0$, donc $1 + \text{e}^{- 3x} > 1 > 0$ et $\left(1 + \text{e}^{- 3x}\right)^2 > 0$ : tous les termes de la dérivée sont supérieurs à zéro ; on a donc $f'(x) > 0$, sur $\R$. La fonction $f$ est strictement croissante sur $\R$.
\item
	\begin{enumerate}
		\item 
Avec $X = - 3x$, on a $\displaystyle\lim_{x \to + \infty} -3x = - \infty$, d'où $\displaystyle\lim_{X \to - \infty} \text{e}^{X} = 0$, d'où $\displaystyle\lim_{x \to + \infty}1 + \text{e}^{-3x} = 1$ et enfin $\displaystyle\lim_{x \to + \infty}f(x) = \dfrac11 = 1$.

\emph{Rem.} la droite d'équation $y = 1 $ est asymptote horizontale à $\mathcal{C}_f$ au voisinage de plus l'infini.
		\item 
On a $\displaystyle\lim_{x \to -\infty}\text{e}^{-3x} = + \infty$, puis $\displaystyle\lim_{x \to -\infty}1 + \text{e}^{-3x}$ et enfin $\displaystyle\lim_{x \to -\infty}\dfrac{1}{1 + \text{e}^{-3x}} = 0$.

\emph{Rem.} la droite d'équation $y = 0$ (l'axe des abscisses) est asymptote horizontale à $\mathcal{C}_f$ au voisinage de moins l'infini.
	\end{enumerate}
\item 
On a $f(x) = 0,99 \iff \dfrac{1}{1 + \text{e}^{- 3x}} = 0,99 \iff 1 = 0,99\left(1 + \text{e}^{- 3x}\right) \iff 1 = 0,99 + 0,99\text{e}^{- 3x} \iff 0,01 = 0,99\text{e}^{- 3x} \iff \dfrac{0,01}{0,99} = \text{e}^{- 3x}\iff \dfrac{1}{99} = \text{e}^{- 3x}$.

Par croissance de la fonction logarithme népérien :

$\ln\left(\dfrac{1}{99}\right) = \ln \left(\text{e}^{- 3x}\right)\iff \ln 1 - \ln 99 = - 3x \iff - \ln 99 = - 3x \iff x = \dfrac{\ln 99}{3}$ \quad (valeur approchée à la calculatrice 1,532).
\end{enumerate}

\bigskip

\textbf{Partie C : Tangente et convexité}

\medskip

\begin{enumerate}
\item Déterminer par le calcul une équation de la tangente $\mathcal{T}$ à la courbe $\mathcal{C}_f$ au point d'abscisse $0$.

On a $M(x~;~y) \in \mathcal{T} \iff y - f(0) = f'(0)(x - 0)$.

$\bullet~~f(0) = \dfrac{1}{1 + 1} = \dfrac{1}{2}$ ;

$\bullet~~f'(0) = \dfrac{3}{(1 + 1)^2} = \dfrac{3}{4}$. Donc :

$M(x~;~y) \in \mathcal{T} \iff y - \dfrac{1}{2} = \dfrac{3}{4}(x - 0) \iff y = \dfrac{3}{4}x + \dfrac{1}{2}$.






\[f''(x) = \dfrac{9\text{e}^{-3x} \left(\text{e}^{-3x} - 1\right)}{\left(1 + \text{e}^{-3x} \right)^3}.\]

\item 
Comme $9\text{e}^{-3x} > 0$ et $\left(1 + \text{e}^{-3x} \right)^3> 0$, le signe de $f''(x)$ est celui de $\text{e}^{-3x} - 1$ :

$\bullet~~$$\text{e}^{-3x} - 1 = 0 \iff \text{e}^{-3x} = 1 \iff - 3x = \ln 1 = 0 \iff \iff x = 0$ ;

$\bullet~~$$\text{e}^{-3x} - 1 > 0 \iff \text{e}^{-3x} > 1 \iff - 3x > \ln 1 = 0 \iff \iff x < 0$ ;

$\bullet~~$$\text{e}^{-3x} - 1 < 0 \iff \text{e}^{-3x} < 1 \iff - 3x < \ln 1 = 0 \iff \iff x > 0$ ;
\item 
	\begin{enumerate}
		\item 
La question précédente a montré que la dérivée seconde est positive sur $]- \infty~;~0]$ : la fonction $f$ est convexe sur cet intervalle.
		\item 
La dérivée seconde s'annule en $x = 0$ en changeant de signe, donc le point A de la courbe représentative d'abscisse 0 est le point d'inflexion de la courbe $\mathcal{C}_f$.
\item  
En chacun de ses points la courbe $\mathcal{C}_f$ est convexe sur $]- \infty~;~0]$ donc au dessus de toutes les tangentes et en  particulier en A la courbe est au dessus de la tangente en $x = 0$, donc en A, donc au dessus de $\mathcal{T}$.

De même en chacun de ses points la courbe $\mathcal{C}_f$ est concave sur $]0~;~+\infty]$ donc au dessous de toutes les tangentes et en  particulier en A la courbe est au dessous de la tangente en $x = 0$, donc en A, donc au dessus de $\mathcal{T}$.

	\end{enumerate}
\end{enumerate}

\medskip
